\unitopen{07}
\begin{candos}
\item podać godzinę oficjalnie (24 h) i potocznie (\de{halb drei}, \de{Viertel nach}, \de{Viertel vor});
\item nazwać dni tygodnia i miesiące oraz podać datę (\de{am dritten Mai});
\item użyć \de{am}, \de{um}, \de{im} i umówić termin;
\item zbudować zdanie z czasem na pierwszym miejscu: \de{Am Montag habe ich ein Meeting.}
\end{candos}

\section{Słownictwo}
\vocabtable{07}

\section{Godzina}
\textbf{Oficjalnie} (rozkłady, e-maile, telefony) czyta się godzinę i minuty w systemie 24-godzinnym:
\de{14:30} = \de{vierzehn Uhr dreißig}. \textbf{Potocznie} dzieli się zegar na kwadranse:

\begin{center}\small
\begin{tabular}{@{}lll@{}}
\toprule
\textsf{\textbf{Zegar}} & \textsf{\textbf{Oficjalnie}} & \textsf{\textbf{Potocznie}}\\
\midrule
08:00 & \de{acht Uhr} & \de{acht Uhr}\\
08:15 & \de{acht Uhr fünfzehn} & \de{Viertel nach acht}\\
08:30 & \de{acht Uhr dreißig} & \de{halb neun}\\
08:45 & \de{acht Uhr fünfundvierzig} & \de{Viertel vor neun}\\
08:10 & \de{acht Uhr zehn} & \de{zehn nach acht}\\
08:50 & \de{acht Uhr fünfzig} & \de{zehn vor neun}\\
\bottomrule
\end{tabular}
\end{center}
\de{halb neun} = pół \emph{do} dziewiątej, czyli 8:30: tak samo jak w polskim „wpół do dziewiątej”.
W niektórych regionach mówi się też \de{Viertel neun} (8:15) lub \de{dreiviertel neun} (8:45); to
regionalizmy, które zrozumiesz, ale sam używaj \de{Viertel nach/vor}.

\section{Dni, miesiące, daty}
Dzień tygodnia i miesiąc są rodzaju męskiego (\artm{der}). Przyimki: \de{am} + dzień (\de{am Montag}), \de{um} + godzina
(\de{um acht Uhr}), \de{im} + miesiąc (\de{im Mai}). \de{im} + pora roku: \de{im Sommer}. \de{am Wochenende}, \de{in der Woche}.

Daty to \emph{liczebniki porządkowe}: do 19. dodaj \de{-te}, od 20. \de{-ste}: \de{der erste} (1.), \de{zweite}, \de{dritte},
\de{vierte}, \de{siebte}, \de{achte}, \de{zwanzigste}, \de{dreißigste}. Z \de{am}: \de{am dritten Mai}, \de{am ersten Januar},
\de{am zwanzigsten Juni}. Zapis: \de{3. Mai}, \de{03.05.2026}. Pytanie: \de{Der Wievielte ist heute?} \de{Den Wievielten haben wir heute?}

\begin{zasada}[Czas pierwszy, czasownik drugi]
Gdy zdanie zaczynasz od okolicznika czasu, podmiot idzie za czasownik: \de{Am Montag \textbf{habe} ich
ein Meeting.} Nie: \de{Am Montag ich habe ein Meeting.} Kolejność w środku zdania: czas, powód, miejsce
(\emph{Te-Ka-Mo-Lo}): \de{Ich fahre morgen mit dem Zug nach Zürich.}
\end{zasada}

\begin{falle}
Po polsku mówisz „o ósmej”, więc \emph{um} wydaje się oczywiste, ale „w poniedziałek” to \de{am Montag}, a nie
\de{in Montag}. Drugi błąd to \de{Wie viel Uhr}: \de{Wie viel Uhr ist es?} jest poprawne, a
\de{Um wie viel Uhr beginnt das Meeting?} to pytanie o \emph{moment}, na które odpowiadasz \de{um}.
Trzeci: \de{morgen} jako przysłówek („jutro”) pisze się małą literą, a \de{der Morgen} (ranek) wielką.
\end{falle}

\begin{szwajcaria}
W Szwajcarii pisze się dzień przed miesiącem (\de{31.12.2026}), a zegar pokazuje 24 h, także na dworcu (SBB).
W biurze punktualność jest normą: \de{pünktlich} znaczy dosłownie na czas, a spóźnienie wymaga
przeprosin i krótkiego uprzedzenia. Posiłki mają dialektalne nazwy: \de{Znüni} (przekąska ok.\ 9 h),
\de{Zmittag}, \de{Zvieri} (przekąska po południu), \de{Znacht}; zrozumiesz je od B1.
\end{szwajcaria}

\begin{praca}
\de{Das Daily ist um halb zehn.} \de{Das Planning ist am Dienstag um vierzehn Uhr.} \de{Kannst du den Termin
auf Donnerstag verschieben?} \de{Ich habe heute um 15 Uhr einen Termin.} Trzy najczęstsze czasowniki
kalendarza: \de{verschieben}, \de{absagen}, \de{vereinbaren} (zapamiętaj jako zwroty).
\end{praca}

\section{Dialogi}
\textbf{Dialog 1 --- umawianie spotkania}
\begin{dialog}
\turn{A}{Haben Sie am Dienstag Zeit?}{Czy ma Pan we wtorek czas?}
\turn{B}{Am Dienstag habe ich um zehn Uhr einen Termin. Und am Mittwoch?}{We wtorek o dziesiątej mam spotkanie. A w środę?}
\turn{A}{Am Mittwoch passt es mir gut. Um wie viel Uhr?}{W środę mi pasuje. O której godzinie?}
\turn{B}{Um halb drei. Ist das okay?}{O wpół do trzeciej. Czy to w porządku?}
\turn{A}{Ja, perfekt. Dann sehen wir uns am Mittwoch, dem vierten Juni, um vierzehn Uhr dreißig.}{Tak, idealnie. Zatem widzimy się w środę, czwartego czerwca, o czternastej trzydzieści.}
\end{dialog}

\section{Ćwiczenia}
\begin{cwiczenia}
\exitem{Podaj godzinę oficjalnie: 14:30.}{\de{vierzehn Uhr dreißig}}
\exitem{Podaj potocznie: 08:15.}{\de{Viertel nach acht}}
\exitem{Podaj potocznie: 17:45.}{\de{Viertel vor sechs}}
\exitem{Podaj potocznie: 09:30.}{\de{halb zehn}}
\exitem{Uzupełnij przyimek: \de{Das Meeting ist \dots{} Montag.}}{\de{am}}
\exitem{Uzupełnij przyimek: \de{Wir beginnen \dots{} zehn Uhr.}}{\de{um}}
\exitem{Uzupełnij przyimek: \de{Ich habe Urlaub \dots{} August.}}{\de{im}}
\exitem{Podaj datę słowami: 3. Mai (z \de{am}).}{\de{am dritten Mai}}
\exitem{Podaj datę słowami: 20. Juni (z \de{am}).}{\de{am zwanzigsten Juni}}
\exitem{Ułóż zdanie, zaczynając od \de{Heute}: (arbeite / ich / im Büro).}{\de{Heute arbeite ich im Büro.}}
\exitem{Popraw: \de{Am Dienstag ich habe einen Termin.}}{\de{Am Dienstag habe ich einen Termin.}}
\exitem{Przetłumacz: „Spotkanie jest we wtorek o czternastej.”}{\de{Das Meeting ist am Dienstag um vierzehn Uhr.}}
\exitem{Zapytaj, o której zaczyna się spotkanie (formalnie).}{\de{Um wie viel Uhr beginnt das Meeting?}}
\end{cwiczenia}

\weekend{Opisz po niemiecku swój typowy tydzień pracy: trzy stałe terminy z dniem i godziną (\de{Montags
um \dots{}}, \de{Am Freitag um \dots{}}). Potem umów się z AI na spotkanie: zaproponuj termin, zaakceptuj
albo przełóż.}
